\section{Multicategories}
\label{sec:multicat}

A \emph{multicategory} $\UU$ consists of:
\begin{itemize}
\item A collection (set, class, etc.) $\Obj \UU$ of \emph{objects}
  $a,b,c,\ldots$.
\item A collection (set, class, etc.) $\UU(a_1, \ldots, a_n; b)$ of
  \emph{multi-morphisms}, for every (possibly empty) list of
  \emph{source} objects $a_1, \ldots, a_n$ and \emph{target} object
  $b$, in which case we write $f : (a_1, \ldots, a_n) \to b$ and
  depict this multi-morphisms like so:
  \[
  \quiver{multi-morphism}
  \] The
  \emph{morphisms} are those multi-morphisms whose source is of length
  $1$.
\item An \emph{identity} morphism $\id[a] : (a) \to a$ for each object $a$.

  We'll depict morphisms by single-source arrows and identity
  morphisms by elongated equality symbols, as in ordinary commutative
  diagrams.

\item The \emph{composition} operation, defined for every
  multi-morphism $g : (b_1, \ldots, b_n) \to c$ and compatible list of
  multi-morphisms $f_i : (a_{i1},\ldots, a_{im_i}) \to b_i$, for each
  $1 \leq i \leq n$, yielding a multi-morphism $g \compose (f_1,
  \ldots, f_n) : (a_{11}, \ldots, a_{1m_1}, \ldots, a_{n1}, \ldots,
  a_{nm_n}) \to c$. We depict the composed morphism by pasting the
  shared objects:
  \[
  \quiver{multi composition}
  \]
\end{itemize}
such that:
\begin{itemize}
\item identities are neutral, i.e., for all $f : (a_1, \ldots, a_n) \to b$:
  \[
  f \compose (\id[a_1], \ldots, \id[a_n]) = f = \id[b] \compose f
  \]
\item composition is associative:
  \begin{multline*}
  h \compose (g_1 \compose (f_{11}, \ldots, f_{1m_1}), \ldots, g_n\compose (f_{n1}, \ldots, f_{nm_n})) \\=
  (h \compose (g_1, \ldots, g_n)) \compose (f_{11}, \ldots, f_{nm_1}, \ldots, f_{n1}, \ldots, f_{nm_n})
  \end{multline*}
  Associativity means that, like for ordinary commutative diagrams, we
  may decompose a multi-morphism diagram in any several equivalent
  ways. For example, we depict the multi morphism of the associativity
  equation like so:
  \[
  \quiver{multi composition associativity}
  \]
\end{itemize}
The paradigmatic example is multi-linear maps:
% \begin{example}
%   Given a field $\Field$, we define the following multicategory
%   $\Multi\Vect_{\Field}$:
%   \begin{itemize}
%   \item Objects are vector spaces $U, V, W, \ldots$ over $\Field$.
%   \item Multi-morphisms $T : (U_1, \ldots, U_n) \to V$ are the
%     \emph{multi-linear} forms/transformations, i.e., the functions $T
%     : U_1 \times \cdots \times U_n \to V$ that are linear in each
%     component:
%     \begin{multline*}
%     T(u_1, \ldots, u_{i-1}, \lambda u_i + \lambda' u'_i, u_{i+1}, \ldots, u_n)
%     \\=
%     \lambda T(u_1, \ldots, u_{i-1}, u_i, u_{i+1}, \ldots, u_n)
%     +
%     \lambda'T(u_1, \ldots, u_{i-1}, u'_i, u_{i+1}, \ldots, u_n)
%     \end{multline*}
%   \end{itemize}
%   For example, the inner product over the field of reals is a
%   bi-linear transformation.

%   We compose multi-linear functions in the evident way:
%   \begin{multline*}
%     g \compose (f_1, \ldots, f_n) (u_{11}, \ldots, u_{1m_1}, \ldots, u_{n1}, \ldots, u_{nm_n})
%     \\:=
%   g(f_1(u_{11}, \ldots, u_{1m_1}), \ldots, f_n(u_{n1}, \ldots, u_{nm_n}))
%   \end{multline*}
%   and the morphisms in $\Multi\Vect_{\Field}$ are the linear
%   transformations, including the identity functions as identities.
% \end{example}

\begin{example}\examplelabel{multi-cat from chosen products}
  Let $\pair{\C}{\prod}$ be a category with specified finite
  products. We equip it with a multicategory structure as follows:
  \begin{itemize}
  \item The objects are the same, $\Obj{\pair\C\prod} := \Obj\C$.
  \item A multi-morphism $f : (a_1, \ldots, a_n) \to b$ is any
    morphism $f : \prod_{i = 1}^n a_i \to b$.
  \item The identities rely on the canonical isomorphism:
    \[
    \id[a] := \prod_{i=1}^1 a \isomorphic a
    \]
  \item Similarly, composition relies on a canonical
    isomorphism. Letting:
    \begin{itemize}
    \item $[n] \definedby \set{1, \ldots, n}$ be the
    enumeration of the index set;
  \item $p := \sum_{i = 1}^n m_i$ the total number of objects in the source of the comosition morphism; and
  \item $\phi : [p] \xto{\isomorphic} \coprod_{i=1}^n[m_i]$ be an
    enumeration of the nested set of indices that's compatible with
    the source list of the composition;
    \end{itemize}
    define:
    \[
    g \compose (f_1, \ldots, f_n) : \prod_{\ell = 1}^{p}a_{\phi \ell}
    \xto{\isomorphic} \prod_{i = 1}^n\prod_{j = 1}^{m_i} a_{i,j}
    \xto{\prod_{i}f_i} \prod_{i=1}^n b_i \xto{g} c
    \]
  \item Tedious verification of the axioms.
  \end{itemize}
\end{example}
Here is a slightly different way to construct a similar multicategory,
but this time without postulating a choice of finite products. Not
sure it's a good way to do it.
\begin{example}\examplelabel{fin prod multi cat}
Let $\C$ be a category that merely has finite products.
\begin{itemize}
\item Objects are the same.
\item A multi-morphism $f : (a_1, \ldots, a_n) \to b$ consists of a
  chosen product $\Dom f$ for $\pi_{i = 1}^na_i$ together with a
  morphism $f : \Dom f \to b$. For identities $\id[a]$ we can then
  take the chosen product $\prod_{i=1}^1$ to be $a$ itself together
  with the identity, and for composition we can take
  $\prod_{i=1}^n\prod_{j=1}^{m_i} a_i$, although here we need to
  employ some choice for the outer product.
\end{itemize}
\end{example}
Generalising the example $\pair\C\prod$:
% \begin{example}
%   Let $\C = \seq{\carrier\C,(\otimes),I,a,l,r}$ be a monoidal
%   category. We equip it with a multi-category structure:
%   \begin{itemize}
%   \item The objects are the same $\Obj\C := \Obj{\carrier\C}$.
%   \item A multi-morphism $f : (a_1, \ldots, a_n) \to b$ is a morphism
%     $f : a_1 \otimes \cdots \otimes a_n \to b$ (say right-nesting the
%     monoidal product).
%   \end{itemize}
%   with the remainder following the same principles as the finite
%   product example.
% \end{example}

A \emph{multi-functor} $F : \UU[1] \to \UU[2]$ consists of two mappings:
\begin{itemize}
\item $F : \Obj{\UU[1]} \to \Obj{\UU[2]}$ sending $\UU[1]$ objects to $\UU[2]$ objects; and
\item $F : \UU[1](a_1, \ldots, a_n; b) \to \UU[2](Fa_1, \ldots, Fa_n;
  Fb)$ sending a multi-morphism $f : (a_1, \ldots, a_n) \to b$ to a
  multi-morphism $Ff : (Fa_1, \ldots, Fa_n) \to Fb$.
\end{itemize}
such that $F$ preserves:
\begin{itemize}
\item identities: $F\id[a] = \id[Fa] : (Fa) \to a$; and
\item composition: $F(g \compose (f_1, \ldots, f_n)) = F g \compose
  (Ff_1, \ldots Ff_n)$.
\end{itemize}
As usual, multi-functors are defined to preserve the structure of
diagrams with multi-morphisms.
\begin{example}
  The identity multi-functor $\Id : \UU \to \UU$ sends every object
  and multi-morphism to themselves. The composition of multi-functors
  is a multi-functor.
\end{example}

A \emph{multi-natural transformation} $\quiver{multi-natural
  transformation type}$ between multi-functors consists of:
\begin{itemize}
\item a morphism $\alpha_a : Fa \to Ga$ for every object $a \in \Obj\UU$
\end{itemize}
such that:
\begin{itemize}
\item
  \[
  \quiver{multi-naturality}
  \]
\end{itemize}

We can define horizontal and vertical composition of multi-natural
transformations analogously to their category theoretic way:
\[
(\beta\vcompose\alpha)_a := \beta_a \compose (\alpha_a) \qquad
(\beta\hcompose\alpha)_a := \beta_{G_1a}\compose (F_2\alpha_a)
                          =G_2\alpha_a\compose (\beta_{F_1a})
\]
and they are multi-natural for analogous reasons. As a consequence we
have a $2$-category $\MultiCAT$ where:
\begin{itemize}
\item $0$-cells are multi-categories;
\item $1$-cells are multi-functors; and
\item $2$-cells are multi-natural transformations
\item with multi-functor composition and identities:
  \[
  \Id[{\CC}] a := a \qquad \Id[{\CC}]f := f
  \]
  vertical, horizontal composition of multi-natural transformations,
  with identities:
  \[
  (\id[{F}])_a := \id[{Fa}] : (a) \to a
  \]

\end{itemize}

For example, we have a $2$-functor structure $\ordinary- : \MultiCAT \to
\CAT$:
\begin{itemize}
\item sending each multicategory $\UU$ to its (\textbf{o}rdinary)
  \emph{category of morphisms}:
  \begin{itemize}
  \item $\Obj {\ordinary\UU} := \Obj \UU$;
  \item $\ordinary\UU(a,b) := \UU(a;b)$, i.e., the morphisms in
    $\ordinary\UU$ are the multi-morphisms with a length-$1$ source.
  \item identities and composition are given by identity morphisms and
    composition in the multi-category.
  \end{itemize}
\item sending each multi-functor $F : \UU[1] \to \UU[2]$ to the
  functor:
  \begin{itemize}
  \item sending $a \in \Obj{\UU[1]}$ to $Fa \in \Obj{\UU[2]}$;
  \item sending $f : (a) \to b$ to $\ordinary Ff := Ff : (Fa) \to Fb$.
  \end{itemize}
\item sending each multi-natural transformation $\quiver{multi-natural
  transformation type}$ to $(\ordinary\alpha)_a := \alpha_a$.
\end{itemize}
\begin{proposition}
  The above assignment is a $2$-functor $\ordinary- : \MultiCAT \to
  \CAT$.

  It is $2$-faithful in the following sense: given two parallel
  multi-natural transformations $\quiver{parallel multi-natural}$, we have
  \(
  \ordinary \alpha = \ordinary \beta \implies \alpha = \beta
  \) and so $\ordinary-$ reflects adjunctions.
\end{proposition}
\begin{proof}
  To show that the $0$-cell mapping is well-defined, we need to show
  we have a category. Composition of morphisms $a \xto f b \xto g c$
  in $\ordinary\UU$ is their composition in the multi-category, i.e.,
  given by \( g \compose f := g \compose (f) \). The multi-category
  axioms then specialise to:
  \[
  \id \compose (f) = f = f \compose (\id)
  \qquad
  h \compose (g \compose (f)) = (h \compose (g)) \compose (f)
  \]
  and so we have a category.

  To show that the $1$-cell mapping is functorial, note that the
  construction type-checks, hence we have a functor structure, and
  then immediately validate functoriality:
  \[
  \ordinary F (\id) = \id
  \qquad
  \ordinary F(g \compose (f)) = F g \compose (F f)
  \]

  For $2$-cells, the resulting transformation $\ordinary\alpha :
  \ordinary F \to \ordinary G$ is natural since the multi-naturality
  condition, when applied to a morphism, recovers the naturality
  condition for the ordinary transformation.

  The $2$-faithfulness follows by observing that the data for the
  multi-natural transformation is preserved by
  $\ordinary-$. Explicitly, if $\ordinary\alpha = \ordinary\beta$,
  then for each $a \in \Obj\UU$, we have $\alpha_a = \beta_a$, and
  that means that $\alpha = \beta$.
\end{proof}
